%! Author = sriadityavedantam
%! Date = 4/8/26

\section{Singularities}

Suppose $f : X\to Y$ is a regular map, where $X,Y$ are quasi-projective varieties.
The question that arises is: if $X$ is irreducible, are the fibers of $f$ irreducible, and if $X$ is smooth, are the fibers smooth?
Generally this is not true.
For example,
\begin{align*}
    \aa^1 &\overset{f}{\longrightarrow} \aa^1 \\
    x &\longmapsto x^2
\end{align*}
where $f^{-1}(y)$ is not irreducible for general $y$.

\begin{theorem*}[First Bertini's Theorem]{
    Assume $X,Y$ are irreducible and $f(X)$ is dense in $Y$, that is, $f$ is dominant.
    Suppose $X$ remains irreducible over the algebraic closure $\overline{\c(Y)}$.
    Then there exists an open dense subset $U\subseteq Y$ such that for all $y\in U$, the fiber $f^{-1}(y)$ is irreducible.
}
\end{theorem*}

If we have the previous example $\aa^1\to \aa^1$, then $\c[Y] = \c[t]$, so $\c(Y) = \c(t)$, which is not algebraically closed.
Since we do not have $t^{1/2}$, and even if we adjoin one root, we still will not have all $t^{1/n}$.

\begin{theorem*}{
    \[
        \overline{\c(t)} = \bigcup_{n\geq 1}\c\bigl(t^{1/n}\bigr),
    \]
    provided $\char 0$.
}
\end{theorem*}

How can we view $X$ as an irreducible variety over $\overline{\c(Y)}$?
Assume $X,Y$ are affine.
Then $f$ induces $\c[Y]\to \c[X]$ by $\varphi \mapsto \varphi\circ f$.
Thus $\c[X]$ is a $\c[Y]$-algebra.
Note that an algebra is a vector space equipped with a bilinear product.
We have
\[
    \c[Y]\subseteq \c(Y)\subseteq \overline{\c(Y)}.
\]
Then consider
\[
    \c[X]\otimes_{\c[Y]}\overline{\c(Y)}.
\]

\begin{definition}[Tensor Product]
    An element of $A\otimes_C B$ satisfies the relation
    \[
        ca\otimes b = a\otimes cb.
    \]
\end{definition}

Thus if we let
\[
    \tilde{X} \coloneqq \spec\left(\c[X]\otimes_{\c[Y]}\overline{\c(Y)}\right),
\]
then $\tilde{X}$ is affine.
In the statement of First Bertini's Theorem, we have that $\tilde{X}$ is irreducible if and only if $X$ is irreducible over $\overline{\c(Y)}$.
Consider
\[
    \tilde{X} = \spec\left(\c[X]\otimes_{\c[t]}\overline{\c(t)}\right).
\]
We claim that $\tilde{X}$ is not irreducible.
Since
\[
    \c[x]\cong \c[t][x]/(x^2-t)
\]
as a $\c[t]$-algebra, we get
\[
    \c[t][x]/(x^2-t)\otimes_{\c[t]}\overline{\c(t)}
    \cong
    \overline{\c(t)}[x]/(x^2-t).
\]
We can now factor over $\overline{\c(t)}$:
\[
    \overline{\c(t)}[x]/(x^2-t)
    =
    \overline{\c(t)}[x]/\left(x-\sqrt{t}\right)\left(x+\sqrt{t}\right).
\]
Thus
\[
    \tilde{X}
    =
    \spec\left(\overline{\c(t)}[x]/\left(x-\sqrt{t}\right)\left(x+\sqrt{t}\right)\right)
\]
has two points, hence is not irreducible.
The First Bertini Theorem ensures irreducibility over points $U\subseteq Y$ open and dense, but not all of $Y$.

\begin{theorem*}[Second Bertini's Theorem]{
    Given $f : X\to Y$ with $X,Y$ quasi-projective, $f(X)$ dense in $Y$, and $X$ nonsingular, there exists an open dense subset $U\subseteq Y$ such that for all $y\in U$, the fiber $f^{-1}(y)$ is smooth.
}
\end{theorem*}

We can think of this theorem as the algebro-geometric analog of Sard's Theorem from differential geometry.

\begin{theorem*}[Sard's Theorem]{
    Let $f : M\to N$ be a smooth map of manifolds.
    Then the set of critical values of $f$ has measure zero in $N$.
}
\end{theorem*}

\begin{definition}[Critical Points]
    A point $x\in M$ is a critical point if and only if
    \[
        d_{x}f : T_{x}M\to T_{f(x)}N
    \]
    is not surjective.
\end{definition}